% !TeX root = main.tex
% ============================================================
%  Tactical Formation — Preamble
%  Compile (XeLaTeX, mặc định):
%    xelatex main && bibtex main && xelatex main && xelatex main
%  Để dùng pdfLaTeX: comment block "XeLaTeX" và bỏ comment block "pdfLaTeX".
% ============================================================

\documentclass[12pt, a4paper]{article}

% ── Encoding & Language (tự nhận engine) ─────────────────────
% Tự phát hiện trình biên dịch: chạy được với CẢ pdfLaTeX (mặc định
% Overleaf) lẫn XeLaTeX/LuaLaTeX. Không cần đổi gì khi chuyển engine.
\usepackage{iftex}
\ifPDFTeX
  % pdfLaTeX: tiếng Việt qua T5 + babel
  \usepackage[utf8]{inputenc}
  \usepackage[T5]{fontenc}
  \usepackage{lmodern}
  \usepackage[english,main=vietnamese]{babel}
\else
  % XeLaTeX / LuaLaTeX: tiếng Việt qua fontspec + polyglossia
  \usepackage{fontspec}
  \usepackage{polyglossia}
  \setdefaultlanguage{vietnamese}
  \setotherlanguage{english}
  % Font Latin Modern có đủ dấu tiếng Việt; nếu hệ không có, thay
  % bằng `\setmainfont{TeX Gyre Termes}` hoặc `\setmainfont{Noto Serif}`.
  \setmainfont{Latin Modern Roman}
\fi

% ── Page Layout ──────────────────────────────────────────────
\usepackage[top=2.5cm, bottom=2.5cm, left=2.5cm, right=2.5cm]{geometry}

% ── Math ─────────────────────────────────────────────────────
\usepackage{amsmath, amssymb}

% ── Tables ───────────────────────────────────────────────────
\usepackage{booktabs}
\usepackage{tabularx}
\usepackage{array}

% ── Graphics & Diagrams ──────────────────────────────────────
\usepackage{tikz}
\usetikzlibrary{arrows.meta, positioning, shapes.geometric}
\usepackage{graphicx}
\usepackage{float}

% ── Colors ───────────────────────────────────────────────────
\usepackage{xcolor}
\definecolor{vertexcolor}{RGB}{50,  100, 200}
\definecolor{centercolor}{RGB}{200, 80,  50}
\definecolor{yellowpiece}{RGB}{240, 200, 0}
\definecolor{greenpiece} {RGB}{50,  180, 80}

% ── Author / Affiliation ─────────────────────────────────────
\usepackage{authblk}

% ── Bibliography (APA 7 qua natbib) ──────────────────────────
\usepackage[round]{natbib}
\bibliographystyle{apalike}

% ── Hyperlinks ───────────────────────────────────────────────
\usepackage{hyperref}
\usepackage{xurl}            % allow line breaks inside long URLs
\setlength{\emergencystretch}{3em}  % avoid overfull lines
\hypersetup{
  colorlinks = true,
  linkcolor  = blue,
  urlcolor   = blue,
  citecolor  = blue,
  pdftitle   = {Cờ Trận (Tactical Formation): Ghi chép và phân tích tính toán một trò chơi cờ dân gian Việt Nam cho 2–4 người chơi},
  pdfauthor  = {Trương Văn Quyết},
  pdfkeywords = {trò chơi dân gian Việt Nam, cờ truyền thống, di sản văn hoá phi vật thể, Paranoid Minimax}
}

% ── Pseudocode ───────────────────────────────────────────────
\usepackage{algorithm}
\usepackage{algpseudocode}

% ── Code Listings ────────────────────────────────────────────
\usepackage{listings}
\lstset{basicstyle=\ttfamily\small, breaklines=true,
        frame=single, captionpos=b}

% ── Misc ─────────────────────────────────────────────────────
\usepackage{enumitem}
\usepackage{caption}
\usepackage{subcaption}
\usepackage{setspace}
\onehalfspacing

% (Các hình bàn cờ nay là ảnh chụp từ trò chơi — xem tools/make_figures.py;
%  không còn dùng macro TikZ vẽ tay.)
